% degraded-mode.tex  --  Transition from normal to degraded retrieval mode
% % degraded-mode.tex  --  Transition from normal to degraded retrieval mode
% % degraded-mode.tex  --  Transition from normal to degraded retrieval mode
% % degraded-mode.tex  --  Transition from normal to degraded retrieval mode
% \input{figures/degraded-mode}
% Requires: tikz, xcolor[dvipsnames], calc, arrows.meta,
%   decorations.pathreplacing loaded in the preamble.

\begin{figure}[htbp]
\centering
\begin{tikzpicture}[
    >=Stealth,
    every node/.style={font=\small},
  ]

  % --- Dimensions ---
  \def\W{10}       % plot width (cm)
  \def\H{4.0}      % plot height (α=1 maps to \H)
  \def\lz{5.0}     % λ₀ position on x-axis
  \def\gam{1.2}    % γ steepness

  % --- Background zones ---
  \fill[MidnightBlue!8] (0, 0) rectangle (3.0, \H+0.4);
  \fill[gray!6]         (3.0, 0) rectangle (7.0, \H+0.4);
  \fill[BrickRed!8]     (7.0, 0) rectangle (\W, \H+0.4);

  % Zone boundaries
  \draw[black!15, thin] (3.0, 0) -- (3.0, \H+0.4);
  \draw[black!15, thin] (7.0, 0) -- (7.0, \H+0.4);

  % --- Zone labels (above plot) ---
  \node[font=\scriptsize\bfseries, text=MidnightBlue!80]
    at (1.5, \H+0.7) {Штатний};
  \node[font=\scriptsize\bfseries, text=gray!60]
    at (5.0, \H+0.7) {Перехідна зона};
  \node[font=\scriptsize\bfseries, text=BrickRed!80]
    at (8.5, \H+0.7) {Деградований};

  % --- Axes ---
  \draw[->, thick, black!70] (-0.2, 0) -- (\W+0.6, 0)
    node[right, font=\small] {$\lambda$\;\scriptsize(подій/год)};
  \draw[->, thick, black!70] (0, -0.2) -- (0, \H+1.0)
    node[above, font=\small] {$\alpha$};

  % Y-axis ticks
  \foreach \yval/\lab in {0/$0$, 2.0/$0{,}5$, 4.0/$1$} {
    \draw[black!50] (-0.12, \yval) -- (0.12, \yval);
    \node[left, font=\scriptsize, black!70] at (-0.18, \yval) {\lab};
  }

  % α = 0.5 reference line
  \draw[black!25, dotted, thin] (0.12, 2.0) -- (\lz, 2.0);

  % λ₀ reference
  \draw[black!40, dashed, thin] (\lz, 0) -- (\lz, 2.0);
  \node[below, font=\scriptsize, black!70] at (\lz, -0.2) {$\lambda_0$};

  % --- Sigmoid curve: α(λ) = 1/(1 + exp(-γ(λ - λ₀))) ---
  \draw[BrickRed!80!black, very thick]
    plot[domain=0:\W, samples=80, smooth]
    (\x, {\H/(1+exp(-\gam*(\x-\lz)))});

  % Curve label
  \node[font=\scriptsize, text=BrickRed!70!black,
        fill=white, inner sep=1.5pt, rounded corners=1pt]
    at (7.6, 2.3) {$\alpha(\lambda)$};

  % --- Scoring annotations inside zones ---
  \node[font=\scriptsize, text=MidnightBlue!80, align=center]
    at (1.5, 2.6)
    {$\mathrm{score} \approx$\\[1pt]$\mathrm{sim}(q,e)$};

  \node[font=\scriptsize, text=BrickRed!80, align=center]
    at (8.5, 2.6)
    {$\mathrm{score} \approx$\\[1pt]$\mathrm{crit}(e)$};

  % --- Digest format annotations (braces below x-axis) ---
  \draw[decorate, decoration={brace, mirror, amplitude=3pt},
        MidnightBlue!60, thick]
    (0, -0.6) -- (3.0, -0.6);
  \node[font=\tiny, text=MidnightBlue!70, below, align=center]
    at (1.5, -0.85) {плоский дайджест\\(перелік фактів)};

  \draw[decorate, decoration={brace, mirror, amplitude=3pt},
        BrickRed!60, thick]
    (7.0, -0.6) -- (\W, -0.6);
  \node[font=\tiny, text=BrickRed!70, below, align=center]
    at (8.5, -0.85) {темпоральна структура\\(динаміка + каузальність)};

\end{tikzpicture}
\caption{%
  Перехід підсистеми пам'яті зі штатного в деградований режим
  як функція частоти подій~$\lambda$.
  При~$\lambda \ll \lambda_0$ ранжування визначається семантичною подібністю
  (штатний режим); при~$\lambda \gg \lambda_0$ домінує тріаж за
  критичністю (деградований режим).
  Сигмоїдний перехід~(\ref{eq:alpha}) забезпечує плавну деградацію
  без~дискретного перемикання.
  Формат дайджесту також адаптується: від плоского переліку фактів
  до темпоральної структури з причинно"=наслідковими зв'язками.%
}
\label{fig:degraded-mode}
\end{figure}

% Requires: tikz, xcolor[dvipsnames], calc, arrows.meta,
%   decorations.pathreplacing loaded in the preamble.

\begin{figure}[htbp]
\centering
\begin{tikzpicture}[
    >=Stealth,
    every node/.style={font=\small},
  ]

  % --- Dimensions ---
  \def\W{10}       % plot width (cm)
  \def\H{4.0}      % plot height (α=1 maps to \H)
  \def\lz{5.0}     % λ₀ position on x-axis
  \def\gam{1.2}    % γ steepness

  % --- Background zones ---
  \fill[MidnightBlue!8] (0, 0) rectangle (3.0, \H+0.4);
  \fill[gray!6]         (3.0, 0) rectangle (7.0, \H+0.4);
  \fill[BrickRed!8]     (7.0, 0) rectangle (\W, \H+0.4);

  % Zone boundaries
  \draw[black!15, thin] (3.0, 0) -- (3.0, \H+0.4);
  \draw[black!15, thin] (7.0, 0) -- (7.0, \H+0.4);

  % --- Zone labels (above plot) ---
  \node[font=\scriptsize\bfseries, text=MidnightBlue!80]
    at (1.5, \H+0.7) {Штатний};
  \node[font=\scriptsize\bfseries, text=gray!60]
    at (5.0, \H+0.7) {Перехідна зона};
  \node[font=\scriptsize\bfseries, text=BrickRed!80]
    at (8.5, \H+0.7) {Деградований};

  % --- Axes ---
  \draw[->, thick, black!70] (-0.2, 0) -- (\W+0.6, 0)
    node[right, font=\small] {$\lambda$\;\scriptsize(подій/год)};
  \draw[->, thick, black!70] (0, -0.2) -- (0, \H+1.0)
    node[above, font=\small] {$\alpha$};

  % Y-axis ticks
  \foreach \yval/\lab in {0/$0$, 2.0/$0{,}5$, 4.0/$1$} {
    \draw[black!50] (-0.12, \yval) -- (0.12, \yval);
    \node[left, font=\scriptsize, black!70] at (-0.18, \yval) {\lab};
  }

  % α = 0.5 reference line
  \draw[black!25, dotted, thin] (0.12, 2.0) -- (\lz, 2.0);

  % λ₀ reference
  \draw[black!40, dashed, thin] (\lz, 0) -- (\lz, 2.0);
  \node[below, font=\scriptsize, black!70] at (\lz, -0.2) {$\lambda_0$};

  % --- Sigmoid curve: α(λ) = 1/(1 + exp(-γ(λ - λ₀))) ---
  \draw[BrickRed!80!black, very thick]
    plot[domain=0:\W, samples=80, smooth]
    (\x, {\H/(1+exp(-\gam*(\x-\lz)))});

  % Curve label
  \node[font=\scriptsize, text=BrickRed!70!black,
        fill=white, inner sep=1.5pt, rounded corners=1pt]
    at (7.6, 2.3) {$\alpha(\lambda)$};

  % --- Scoring annotations inside zones ---
  \node[font=\scriptsize, text=MidnightBlue!80, align=center]
    at (1.5, 2.6)
    {$\mathrm{score} \approx$\\[1pt]$\mathrm{sim}(q,e)$};

  \node[font=\scriptsize, text=BrickRed!80, align=center]
    at (8.5, 2.6)
    {$\mathrm{score} \approx$\\[1pt]$\mathrm{crit}(e)$};

  % --- Digest format annotations (braces below x-axis) ---
  \draw[decorate, decoration={brace, mirror, amplitude=3pt},
        MidnightBlue!60, thick]
    (0, -0.6) -- (3.0, -0.6);
  \node[font=\tiny, text=MidnightBlue!70, below, align=center]
    at (1.5, -0.85) {плоский дайджест\\(перелік фактів)};

  \draw[decorate, decoration={brace, mirror, amplitude=3pt},
        BrickRed!60, thick]
    (7.0, -0.6) -- (\W, -0.6);
  \node[font=\tiny, text=BrickRed!70, below, align=center]
    at (8.5, -0.85) {темпоральна структура\\(динаміка + каузальність)};

\end{tikzpicture}
\caption{%
  Перехід підсистеми пам'яті зі штатного в деградований режим
  як функція частоти подій~$\lambda$.
  При~$\lambda \ll \lambda_0$ ранжування визначається семантичною подібністю
  (штатний режим); при~$\lambda \gg \lambda_0$ домінує тріаж за
  критичністю (деградований режим).
  Сигмоїдний перехід~(\ref{eq:alpha}) забезпечує плавну деградацію
  без~дискретного перемикання.
  Формат дайджесту також адаптується: від плоского переліку фактів
  до темпоральної структури з причинно"=наслідковими зв'язками.%
}
\label{fig:degraded-mode}
\end{figure}

% Requires: tikz, xcolor[dvipsnames], calc, arrows.meta,
%   decorations.pathreplacing loaded in the preamble.

\begin{figure}[htbp]
\centering
\begin{tikzpicture}[
    >=Stealth,
    every node/.style={font=\small},
  ]

  % --- Dimensions ---
  \def\W{10}       % plot width (cm)
  \def\H{4.0}      % plot height (α=1 maps to \H)
  \def\lz{5.0}     % λ₀ position on x-axis
  \def\gam{1.2}    % γ steepness

  % --- Background zones ---
  \fill[MidnightBlue!8] (0, 0) rectangle (3.0, \H+0.4);
  \fill[gray!6]         (3.0, 0) rectangle (7.0, \H+0.4);
  \fill[BrickRed!8]     (7.0, 0) rectangle (\W, \H+0.4);

  % Zone boundaries
  \draw[black!15, thin] (3.0, 0) -- (3.0, \H+0.4);
  \draw[black!15, thin] (7.0, 0) -- (7.0, \H+0.4);

  % --- Zone labels (above plot) ---
  \node[font=\scriptsize\bfseries, text=MidnightBlue!80]
    at (1.5, \H+0.7) {Штатний};
  \node[font=\scriptsize\bfseries, text=gray!60]
    at (5.0, \H+0.7) {Перехідна зона};
  \node[font=\scriptsize\bfseries, text=BrickRed!80]
    at (8.5, \H+0.7) {Деградований};

  % --- Axes ---
  \draw[->, thick, black!70] (-0.2, 0) -- (\W+0.6, 0)
    node[right, font=\small] {$\lambda$\;\scriptsize(подій/год)};
  \draw[->, thick, black!70] (0, -0.2) -- (0, \H+1.0)
    node[above, font=\small] {$\alpha$};

  % Y-axis ticks
  \foreach \yval/\lab in {0/$0$, 2.0/$0{,}5$, 4.0/$1$} {
    \draw[black!50] (-0.12, \yval) -- (0.12, \yval);
    \node[left, font=\scriptsize, black!70] at (-0.18, \yval) {\lab};
  }

  % α = 0.5 reference line
  \draw[black!25, dotted, thin] (0.12, 2.0) -- (\lz, 2.0);

  % λ₀ reference
  \draw[black!40, dashed, thin] (\lz, 0) -- (\lz, 2.0);
  \node[below, font=\scriptsize, black!70] at (\lz, -0.2) {$\lambda_0$};

  % --- Sigmoid curve: α(λ) = 1/(1 + exp(-γ(λ - λ₀))) ---
  \draw[BrickRed!80!black, very thick]
    plot[domain=0:\W, samples=80, smooth]
    (\x, {\H/(1+exp(-\gam*(\x-\lz)))});

  % Curve label
  \node[font=\scriptsize, text=BrickRed!70!black,
        fill=white, inner sep=1.5pt, rounded corners=1pt]
    at (7.6, 2.3) {$\alpha(\lambda)$};

  % --- Scoring annotations inside zones ---
  \node[font=\scriptsize, text=MidnightBlue!80, align=center]
    at (1.5, 2.6)
    {$\mathrm{score} \approx$\\[1pt]$\mathrm{sim}(q,e)$};

  \node[font=\scriptsize, text=BrickRed!80, align=center]
    at (8.5, 2.6)
    {$\mathrm{score} \approx$\\[1pt]$\mathrm{crit}(e)$};

  % --- Digest format annotations (braces below x-axis) ---
  \draw[decorate, decoration={brace, mirror, amplitude=3pt},
        MidnightBlue!60, thick]
    (0, -0.6) -- (3.0, -0.6);
  \node[font=\tiny, text=MidnightBlue!70, below, align=center]
    at (1.5, -0.85) {плоский дайджест\\(перелік фактів)};

  \draw[decorate, decoration={brace, mirror, amplitude=3pt},
        BrickRed!60, thick]
    (7.0, -0.6) -- (\W, -0.6);
  \node[font=\tiny, text=BrickRed!70, below, align=center]
    at (8.5, -0.85) {темпоральна структура\\(динаміка + каузальність)};

\end{tikzpicture}
\caption{%
  Перехід підсистеми пам'яті зі штатного в деградований режим
  як функція частоти подій~$\lambda$.
  При~$\lambda \ll \lambda_0$ ранжування визначається семантичною подібністю
  (штатний режим); при~$\lambda \gg \lambda_0$ домінує тріаж за
  критичністю (деградований режим).
  Сигмоїдний перехід~(\ref{eq:alpha}) забезпечує плавну деградацію
  без~дискретного перемикання.
  Формат дайджесту також адаптується: від плоского переліку фактів
  до темпоральної структури з причинно"=наслідковими зв'язками.%
}
\label{fig:degraded-mode}
\end{figure}

% Requires: tikz, xcolor[dvipsnames], calc, arrows.meta,
%   decorations.pathreplacing loaded in the preamble.

\begin{figure}[htbp]
\centering
\begin{tikzpicture}[
    >=Stealth,
    every node/.style={font=\small},
  ]

  % --- Dimensions ---
  \def\W{10}       % plot width (cm)
  \def\H{4.0}      % plot height (α=1 maps to \H)
  \def\lz{5.0}     % λ₀ position on x-axis
  \def\gam{1.2}    % γ steepness

  % --- Background zones ---
  \fill[MidnightBlue!8] (0, 0) rectangle (3.0, \H+0.4);
  \fill[gray!6]         (3.0, 0) rectangle (7.0, \H+0.4);
  \fill[BrickRed!8]     (7.0, 0) rectangle (\W, \H+0.4);

  % Zone boundaries
  \draw[black!15, thin] (3.0, 0) -- (3.0, \H+0.4);
  \draw[black!15, thin] (7.0, 0) -- (7.0, \H+0.4);

  % --- Zone labels (above plot) ---
  \node[font=\scriptsize\bfseries, text=MidnightBlue!80]
    at (1.5, \H+0.7) {Штатний};
  \node[font=\scriptsize\bfseries, text=gray!60]
    at (5.0, \H+0.7) {Перехідна зона};
  \node[font=\scriptsize\bfseries, text=BrickRed!80]
    at (8.5, \H+0.7) {Деградований};

  % --- Axes ---
  \draw[->, thick, black!70] (-0.2, 0) -- (\W+0.6, 0)
    node[right, font=\small] {$\lambda$\;\scriptsize(подій/год)};
  \draw[->, thick, black!70] (0, -0.2) -- (0, \H+1.0)
    node[above, font=\small] {$\alpha$};

  % Y-axis ticks
  \foreach \yval/\lab in {0/$0$, 2.0/$0{,}5$, 4.0/$1$} {
    \draw[black!50] (-0.12, \yval) -- (0.12, \yval);
    \node[left, font=\scriptsize, black!70] at (-0.18, \yval) {\lab};
  }

  % α = 0.5 reference line
  \draw[black!25, dotted, thin] (0.12, 2.0) -- (\lz, 2.0);

  % λ₀ reference
  \draw[black!40, dashed, thin] (\lz, 0) -- (\lz, 2.0);
  \node[below, font=\scriptsize, black!70] at (\lz, -0.2) {$\lambda_0$};

  % --- Sigmoid curve: α(λ) = 1/(1 + exp(-γ(λ - λ₀))) ---
  \draw[BrickRed!80!black, very thick]
    plot[domain=0:\W, samples=80, smooth]
    (\x, {\H/(1+exp(-\gam*(\x-\lz)))});

  % Curve label
  \node[font=\scriptsize, text=BrickRed!70!black,
        fill=white, inner sep=1.5pt, rounded corners=1pt]
    at (7.6, 2.3) {$\alpha(\lambda)$};

  % --- Scoring annotations inside zones ---
  \node[font=\scriptsize, text=MidnightBlue!80, align=center]
    at (1.5, 2.6)
    {$\mathrm{score} \approx$\\[1pt]$\mathrm{sim}(q,e)$};

  \node[font=\scriptsize, text=BrickRed!80, align=center]
    at (8.5, 2.6)
    {$\mathrm{score} \approx$\\[1pt]$\mathrm{crit}(e)$};

  % --- Digest format annotations (braces below x-axis) ---
  \draw[decorate, decoration={brace, mirror, amplitude=3pt},
        MidnightBlue!60, thick]
    (0, -0.6) -- (3.0, -0.6);
  \node[font=\tiny, text=MidnightBlue!70, below, align=center]
    at (1.5, -0.85) {плоский дайджест\\(перелік фактів)};

  \draw[decorate, decoration={brace, mirror, amplitude=3pt},
        BrickRed!60, thick]
    (7.0, -0.6) -- (\W, -0.6);
  \node[font=\tiny, text=BrickRed!70, below, align=center]
    at (8.5, -0.85) {темпоральна структура\\(динаміка + каузальність)};

\end{tikzpicture}
\caption{%
  Перехід підсистеми пам'яті зі штатного в деградований режим
  як функція частоти подій~$\lambda$.
  При~$\lambda \ll \lambda_0$ ранжування визначається семантичною подібністю
  (штатний режим); при~$\lambda \gg \lambda_0$ домінує тріаж за
  критичністю (деградований режим).
  Сигмоїдний перехід~(\ref{eq:alpha}) забезпечує плавну деградацію
  без~дискретного перемикання.
  Формат дайджесту також адаптується: від плоского переліку фактів
  до темпоральної структури з причинно"=наслідковими зв'язками.%
}
\label{fig:degraded-mode}
\end{figure}
